\section{Experiments}
\label{sec:experiments}

We ask four questions.
\textbf{(i)}~Does a released policy's proposal bank contain better candidates, and can a generic feature distance find them (\cref{sec:headroom})?
\textbf{(ii)}~Does \method select better than the policy and than matched baselines in closed loop (\cref{sec:main})?
\textbf{(iii)}~Where along the chain from actual future to predicted code is ranking lost (\cref{sec:ladder})?
\textbf{(iv)}~Which design choices matter, and does a predicted code transfer to held-out goals better than a direct scorer (\cref{sec:ablations,sec:heldout})?

\subsection{Setup}
\label{sec:setup}

\paragraph{Arenas and policies.}
\emph{PushT}~\cite{florence2021ibc}: we use the released LeRobot Diffusion Policy~\cite{chi2023diffusion}, which observes two frames, denoises a 16-step action tensor, and executes 8 actions.
An episode succeeds when the block covers more than 95\% of the target within 300 steps.
The unmodified policy succeeds on 65 of 100 evaluation episodes; the model card reports 65.4\%.
\emph{LIBERO-Goal}~\cite{liu2023libero}: ten goals in one kitchen scene, with the released SmolVLA~\cite{shukor2025smolvla} checkpoint executing 10 actions per chunk.
With the unmodified evaluation script it reaches 71.3\% over three seeds, below the 91\% reported; we use the policy as measured and do not tune it.
\TODO{LIBERO branching harness, clone fidelity, and the split into training and held-out goals.}

\paragraph{Proposal banks.}
At every decision the policy draws $K{=}8$ chunks from independently seeded noise; the bank is shared by all methods, and $\bA^{(1)}$ is the policy's own draw.
Simulator clones used for branching are verified to reproduce the live continuation of the episode.

\paragraph{Data and splits.}
On PushT, training uses 800 policy rollouts with all eight proposals branched at every decision.
Separate episodes serve offline development (100), closed-loop development (200), and a sealed confirmation set (400).
The confirmation set is opened once, for the locked configuration trained with three seeds.
Goal images are the final frames of 16 successful episodes outside all splits.

\paragraph{Compared scorers.}
All scorers rank the same bank.
\emph{Policy only} executes $\bA^{(1)}$.
The \emph{oracle} simulates every candidate and picks the highest 8-step progress label; it is privileged and bounds myopic selection.
The \emph{actual-future reader} is a reader of the same size, trained with the same data and loss on the uncompressed actual future; it is not deployable and bounds any predictor of the future.
\emph{Code of the actual future} applies \method's reader to the target encoder's code of the simulated future.
The \emph{direct scorer} $D(C,\bA,g)$ has eight layers and is trained with the same data, loss, and steps.
The \emph{per-step latent world model} is a DINO-WM-style predictor of the patch features of every future step, trained on the same data and read by the actual-future reader.
We also compare with \emph{GPC}~\cite{qi2025gpc}.
\TODO{Per-step latent world model and GPC are not implemented yet.}
Finally, \emph{feature distance} scores the actual final frame by its negative mean squared DINOv2 patch-feature distance to the mean feature map of the goal images, which is DINO-WM's planning cost under perfect prediction.

\paragraph{Metrics.}
The primary metric is closed-loop success.
Differences are paired by episode, with bootstrap confidence intervals (10{,}000 resamples) and exact McNemar tests.
Offline, on development decisions, we report the within-bank Spearman correlation between a score and the 8-step label.
We also report the retained gap: the fraction of the oracle's per-decision label gain over $\bA^{(1)}$ that the scorer's choice obtains.

\subsection{Headroom in the proposal bank}
\label{sec:headroom}

\begin{table}[t]
\centering
\small
\setlength{\tabcolsep}{4pt}
\begin{tabular}{@{}lcc@{}}
\toprule
Selection among $K{=}8$ proposals & Set A & Set B \\
 & \scriptsize(100 ep.) & \scriptsize(400 ep.) \\
\midrule
Policy, unmodified inference & 65 & -- \\
Policy only ($\bA^{(1)}$, our sampler) & 64 & 62.5 \\
Bank medoid (no future information) & 68 & -- \\
Feature distance on the true future & 67 & -- \\
Learned reader on the true future & -- & 71.0 \\
Oracle, 8-step progress (privileged) & 79 & 76.5 \\
\bottomrule
\end{tabular}
\caption{\textbf{PushT headroom}: closed-loop success (\%) when selecting among the policy's own proposals, on two disjoint sets of episodes. The two rows on the true future are not deployable; they bound what a world model that predicted the future perfectly could obtain with that scorer.}
\label{tab:headroom}
\end{table}

\Cref{tab:headroom} shows that the bank has headroom.
On set B, selecting by the privileged oracle raises success by 14 points, 95\% CI $[8.8,19.3]$ (94 episodes solved only by the oracle, 38 only by the policy; McNemar $p{<}10^{-5}$).
On set A the gain is 15 points $[5,25]$, and on a third set of 100 episodes it was 10 points $[-1,22]$.
Per decision the gain is small: a strictly better candidate exists at 54\% of decisions, and the oracle's choice improves coverage by only 0.007 on average.
Swapping a single chunk barely changes the outcome (+2 points $[-3,7.5]$ in a split-seed estimate).
Success therefore comes from being slightly better at many decisions, and a scorer must resolve small differences in progress at every one of them.

Feature distance fails to do so even on the true future.
On set A it gains 3 points $[-8,14]$, and its within-bank Spearman correlation with the label is 0.06 over 1{,}742 decisions.
A per-step latent world model that plans with this cost is thus capped near the policy even under perfect prediction.
A reader trained with the ranking loss of \cref{sec:reader} on the true future does much better.
On set B it gains 8.5 points $[2.8,14.3]$ (McNemar $p{=}0.005$), 61\% of the oracle's gain, with a within-bank Spearman correlation of 0.23.
This learned reader on the true future is the upper bound that \method approaches from compressed and then predicted futures.

\subsection{Closed-loop results}
\label{sec:main}

\begin{table}[t]
\centering
\small
\setlength{\tabcolsep}{3pt}
\begin{tabular}{@{}lcccc@{}}
\toprule
 & PushT & \multicolumn{2}{c}{LIBERO-Goal} & ms per \\
 & & seen & held-out & decision \\
\midrule
Policy only & \result{--} & \result{--} & \result{--} & \result{--} \\
Oracle$^\dagger$ & \result{--} & \result{--} & \result{--} & -- \\
Actual-future reader$^\dagger$ & \result{--} & \result{--} & \result{--} & -- \\
Direct scorer & \result{--} & \result{--} & \result{--} & \result{--} \\
Per-step latent WM & \result{--} & \result{--} & \result{--} & \result{--} \\
GPC~\cite{qi2025gpc} & \result{--} & \result{--} & \result{--} & \result{--} \\
\method (ours) & \result{--} & \result{--} & \result{--} & \result{--} \\
\bottomrule
\end{tabular}
\caption{\textbf{Closed-loop success} (\%) on sealed episodes, averaged over three training seeds, and planning latency per decision, including proposal sampling. $^\dagger$Not deployable: simulates every candidate.}
\label{tab:main}
\end{table}

\result{[Main result: \method vs.\ policy only (primary contrast), vs.\ the direct scorer, and retention relative to the actual-future reader and the oracle.]}

\subsection{Attribution ladder}
\label{sec:ladder}

\begin{table}[t]
\centering
\small
\begin{tabular}{@{}lcc@{}}
\toprule
Scorer input & Spearman & Retained gap \\
\midrule
Actual future (uncompressed) & \result{--} & \result{--} \\
Code of the actual future & \result{--} & \result{--} \\
Predicted code, greedy & \result{--} & \result{--} \\
Predicted code, expected codeword & \result{--} & \result{--} \\
Predicted code, 4 samples & \result{--} & \result{--} \\
Direct scorer (actions) & \result{--} & \result{--} \\
\bottomrule
\end{tabular}
\caption{\textbf{Attribution ladder} on PushT development decisions: within-bank Spearman correlation with the 8-step label and retained gap, with episode-clustered 95\% intervals. A drop from one row to the next locates the loss: compression (row 1 to 2) or prediction (row 2 to 3).}
\label{tab:ladder}
\end{table}

\Cref{tab:ladder} scores the same decisions with the same reader architecture as the input moves from the actual future, to its code, to the code predicted from actions.
\result{[Interpretation; code diagnostics: perplexity, distinct codes per bank, $I(S;\bA\mid C)$ in bits, feature $R^2$.]}

\subsection{Ablations}
\label{sec:ablations}
Each ablation changes one design choice, at the same history, proposal bank, data, and training budget.
\begin{itemize}[leftmargin=*,topsep=2pt,itemsep=1pt]
\item \emph{Conditioning on history.} The target encoder does not see $C$, at the same number of bits. The heads still see $C$.
\item \emph{Segment versus endpoint.} The target encoder sees only the final image. A decision is made once per chunk, but some differences between chunks lie in the path.
\item \emph{Segment versus per-step codes.} A conditional code per time step with the same context and the same total number of bits. If it matches \method, the segment claim fails. \TODO{Not implemented.}
\item \emph{Feature decoder.} $\alpha{=}0$ tests whether reconstruction is what keeps the code useful for goals outside training.
\item \emph{Predictability.} $\lambda{=}0$ versus co-design with $\lambda{>}0$ in~\eqref{eq:source}.
\item \emph{Rate.} $M\in\{8,16,32\}$, which traces the accuracy--latency frontier.
\end{itemize}
\result{[Ablation table.]}

\subsection{Held-out goals}
\label{sec:heldout}
The advantage of predicting a goal-independent code over scoring actions directly should appear on goals that were not seen in training.
On LIBERO-Goal, the reader and the direct scorer are trained with labels for a subset of the ten goals, and both are evaluated on the remaining goals with the same world model and bank.
\result{[Held-out goal result.]}
